\section{Methodology}
\label{sec:methodology}

\subsection{Score function}

Both the user and item representations are sums of feature embeddings
(Equation~\ref{eq:feature-sum}, repeated here for the score derivation):
\begin{equation}
  \pvec_u = \sum_{f \in F_u} \evec_f, \qquad
  \qvec_i = \sum_{g \in F_i} \evec_g .
\end{equation}
The score of a user--item pair is their inner product plus an item bias,
\begin{equation}
  s(u, i) = \Big\langle \sum_{f \in F_u} \pvec_f,\ \sum_{g \in F_i} \qvec_g \Big\rangle
          = \langle \pvec_u, \qvec_i \rangle + b_i ,
  \label{eq:score}
\end{equation}
where $b_i = \sum_{g \in F_i} \beta_g$ is itself a sum of per-feature bias
scalars, mirroring the embedding sum. A per-user bias is deliberately
omitted: it would shift every item's score for that user by the same
constant and therefore cancel out of the pairwise margin used by the loss
below, exactly as the per-user bias does in the sibling BPR project's own
derivation -- it would carry zero gradient under this objective.

\subsection{WARP loss}
\label{sec:warp}

Weighted Approximate-Rank Pairwise loss~\cite{weston2011warp,kula2015}
optimizes a pairwise ranking criterion like BPR, but instead of weighting
every violated pair equally, it weights a positive pair by an \emph{estimate
of the positive's current rank} in the full catalogue, so that correcting a
badly mis-ranked positive contributes far more gradient than nudging one
that is already well-placed.

For a positive pair $(u, i)$, negatives $j$ are sampled uniformly and
rejected until a margin violation is found,
\begin{equation}
  s(u, j) > s(u, i) - 1 ,
\end{equation}
or until $N_{\max}$ trials are exhausted. Let $N$ be the number of trials
taken. Because negatives are sampled uniformly at random, the probability
that a given negative violates the margin is approximately
$\text{rank}(i) / (|\Iset| - 1)$, so the expected number of trials before
the first violation is the reciprocal of that probability; inverting gives
the rank estimate
\begin{equation}
  \text{rank}(i) \approx \left\lfloor \frac{|\Iset| - 1}{N} \right\rfloor .
  \label{eq:rank-estimate}
\end{equation}
Few trials imply many items already violate the margin, i.e.\ $i$ is ranked
low; many trials imply $i$ is already hard to beat, i.e.\ ranked high. The
hinge loss for the triplet is weighted by the margin-ranking loss function
\begin{equation}
  L(\text{rank}(i)) = \sum_{r=1}^{\text{rank}(i)} \frac{1}{r} ,
  \label{eq:warp-weight}
\end{equation}
a harmonic sum that grows (slowly) with rank, so that
\begin{equation}
  \ell(u, i, j) = L(\text{rank}(i)) \cdot \max\!\big(0,\ 1 - s(u,i) + s(u,j)\big) .
  \label{eq:warp-loss}
\end{equation}

\subsection{Gradients and the update rule}

Write $m = 1 - s(u,i) + s(u,j)$ for the margin and $w = L(\text{rank}(i))$
for the WARP weight; $\ell = w \cdot \max(0, m)$ is non-zero, and has a
gradient, only when $m > 0$. For an active triplet,
\begin{align}
  \frac{\partial \ell}{\partial s(u,i)} &= -w , &
  \frac{\partial \ell}{\partial s(u,j)} &= +w .
\end{align}
Since $s(u,i) = \langle \pvec_u, \qvec_i \rangle + b_i$ and every feature
embedding is shared across every entity that carries that feature, the
chain rule distributes the gradient to \emph{every active feature
embedding} of $u$, $i$ and $j$ simultaneously:
\begin{align}
  \pvec_u \text{'s features } f &:\ \evec_f \leftarrow \evec_f +
    \eta \big( w (\qvec_i - \qvec_j) - \lambda \evec_f \big) , \\
  \qvec_i \text{'s features } g &:\ \evec_g \leftarrow \evec_g +
    \eta \big( w \, \pvec_u - \lambda \evec_g \big) , \\
  \qvec_j \text{'s features } g &:\ \evec_g \leftarrow \evec_g +
    \eta \big( {-w} \, \pvec_u - \lambda \evec_g \big) , \\
  b_i \text{'s features } g &:\ \beta_g \leftarrow \beta_g + \eta \big( w - \lambda \beta_g \big) , \\
  b_j \text{'s features } g &:\ \beta_g \leftarrow \beta_g + \eta \big( {-w} - \lambda \beta_g \big) ,
\end{align}
with learning rate $\eta$ and L2 penalty $\lambda$. This is the sense in
which the model is a classical, linear SGD factorization: every update is a
closed-form linear-algebra step on embeddings, with no hidden layer and no
nonlinear activation anywhere in the forward pass.

\paragraph{A stability correction that matters in practice.} A feature
shared by many low-popularity training items (in our pipeline this would
have been \texttt{popularity=0}, had it not been excluded from the
representation per Section~\ref{sec:features}) is sampled as a WARP
negative about as often as any other feature but is rarely part of a
positive, so with only a weak L2 pull its bias term can diverge far faster
than $\lambda$ counteracts it -- in early experiments this drove cold-item
scores to the bottom of the ranking regardless of their brand/ingredient
match. Clipping every touched embedding to a bounded L2 norm and every
touched bias scalar to a bounded magnitude after each update (both
configurable; LightFM's own reference implementation does the same for
embeddings) removes the divergence without changing the gradient direction,
only its magnitude.

\subsection{Baselines}

\paragraph{ID-only MF (BPR).} The same pairwise-SGD family, but the user and
item representations are raw identity rows/columns rather than feature
sums -- structurally identical to the sibling BPR project. A user or item
absent from training has no row or column, so its score collapses to the
item bias term alone (a pure popularity signal) for warm items, and to
\emph{nothing meaningful at all} for a cold item, since that item's bias was
also never trained; its factor and bias are masked out of the rank for every
user, warm or cold.

\paragraph{Most-popular.} Ranks every user by the same global training-count
order. Just like the ID-only baseline, a cold item has a training count of
exactly zero and is therefore never recommended by this baseline either.

\paragraph{Random.} A seeded uniform permutation per user, which calibrates
the metric scale; because its coverage is total and uncorrelated with
quality, high coverage alone is not evidence of a good model.
